\subsection{Limiting oracle risks}
\label{sec:oracle-limits}

The finite-prevalence theorem gives explicit gaps for two signal families. We
next characterize both oracles for arbitrary fixed signals and noise as rare
prevalence decreases. Write $s=(\beta_C^*)^2$, $r=(\beta_R^*)^2$,
$\nu=\sigma^2$, and $v=s/4+r$. The limiting initial rare-subgroup risk is $v/2$.
For $\alpha\in[0,1]$, define
\[
  d(\alpha)=\alpha(2-\alpha),\qquad
  z_\tau(\alpha)=e^{-\tau d(\alpha)},\qquad
  g(\alpha)=\frac{\nu\alpha}{8-\alpha},
\]
and, suppressing the arguments of $z_\tau$ and $g$,
\begin{equation}
  \begin{aligned}
    F_S(\alpha;s,r,\nu,\tau)
      &=\frac12\left[vz_\tau+
        \frac{\nu\alpha}{2-\alpha}(1-z_\tau)\right],\\
    F_U(\alpha;r,\nu,\tau)
      &=\frac12\left[g+rz_\tau+
        \frac{\alpha(\nu+g)}{2-\alpha}(1-z_\tau)\right].
  \end{aligned}
  \label{eq:limiting-objectives}
\end{equation}

\begin{theorem}[Limiting oracle risks]
  \label{thm:oracle-limits}
  In the model of \cref{sec:conference-setting}, fix finite $s,r,\nu\geq0$.
  If $q\to0$ and integer observation caps satisfy
  $M_qb\to\tau\in(0,\infty)$, then
  \[
    R_A^*(M_q)\longrightarrow L_A(\tau)
      :=\min_{0\leq\alpha\leq1}F_A(\alpha;\tau),
      \qquad A\in\{U,S\},
  \]
  with the signal and noise arguments as in \cref{eq:limiting-objectives}.
  More precisely, for every $0<T_0<T_1<\infty$, uniformly over integer caps
  with $T_q=M_qb\in[T_0,T_1]$,
  \begin{equation}
    |R_S^*(M_q)-L_S(T_q)|\leq C_Sb,\qquad
    |R_U^*(M_q)-L_U(T_q)|\leq C_Ub\log(1/b)
    \label{eq:oracle-limit-errors}
  \end{equation}
  for all sufficiently small $q$. The constants depend only on
  $s,r,\nu,T_0,T_1$.
\end{theorem}

The oracle in this theorem still searches every nonnegative constant rate and
every deterministic stop, including prevalence-dependent choices and unstable
rates. Reduction to $\alpha\in[0,1]$ and the full cap is a conclusion about the
limiting optimization. The approximation constants are not optimized numerical
certificates at a specified prevalence; \cref{thm:conference-predictive}
provides that separate guarantee.

\paragraph{Why the limits differ.}

SD's limiting risk depends on the signal only through $v$. U's depends on $r$
but loses its dependence on the initial common signal. At a positive limiting
rate, U learns that signal quickly and retains common-coordinate noise $g$,
which also feeds into the rare error. By taking $b\ll\alpha_q\ll1$, U instead
learns the common signal with vanishing noise cost and leaves the rare effect
unchanged. This gives the endpoint $F_U(0)=r/2$. Literal zero-rate training
retains both signal errors and has limiting risk $v/2$.

For $\nu>0$, U uses this endpoint exactly when $r\tau\leq\nu/16$; above that
threshold, its unique minimizer is positive. SD has a unique positive minimizer
whenever $v,\nu>0$. The scalar equations in \cref{sec:oracle-variational}
determine these rates without a grid search.

\paragraph{Proof outline.}

Let $G$ and $H$ denote the common and rare contributions before the factor
$1/2$ in the limiting risk. Under SD, their sum obeys
$\dot W=-d(\alpha)W+\alpha^2\nu$, starting at $W(0)=v$. Under U, the common
moment rapidly approaches $g(\alpha)$, and the rare moment obeys
$\dot H=-d(\alpha)H+\alpha^2(\nu+g)$, starting at $H(0)=r$. Solving these
equations gives \cref{eq:limiting-objectives}. Passing to the oracle infima
also requires uniform bounds for large rates, early stops, and vanishing rates.
We retain U's common-coordinate transient in those bounds and construct rates
that attain each limiting minimum, including its zero endpoint.
\Cref{sec:oracle-limit-proofs} gives the full argument.

\subsection{A unique signal-allocation crossing}
\label{sec:oracle-allocation}

Hold the initial limiting subgroup risk fixed while moving signal between the
features. With $v>0$, write $r=\theta v$ and $s=4(1-\theta)v$ for
$\theta\in[0,1]$. Thus $\theta$ measures the rare feature's share of $s/4+r$,
rather than its share of the unweighted squared coefficient sum.

\begin{corollary}[Signal allocation]
  \label{cor:oracle-allocation}
  Fix $v,\nu,\tau>0$. Along this allocation path, $L_S$ is constant and $L_U$
  increases strictly with $\theta$. There is a unique
  $\theta_*\in(0,1)$ such that $L_S=L_U$. Unscaled training has lower limiting
  risk for $\theta<\theta_*$, and standardization has lower limiting risk for
  $\theta>\theta_*$.
\end{corollary}

\begin{proof}
  The SD objective depends only on $v$. Each U objective has slope
  $vz_\tau(\alpha)/2\in[ve^{-\tau}/2,v/2]$ in $\theta$, so its minimum is
  continuous and strictly increasing. At $\theta=0$, U has zero risk and SD
  has positive risk. At $\theta=1$, a positive U minimizer incurs strictly
  more noise than the same SD rate. If U instead minimizes at zero, a small
  positive SD rate improves on $v/2$. The endpoint signs give the unique crossing.
\end{proof}

\Cref{fig:intro-signal-crossing} illustrates this corollary. At $v=\nu=1$
and $\tau=2$, the crossing is numerically $\theta_*\approx0.8377$.
\Cref{sec:oracle-variational} gives a variational formula and an equivalent
boundary obtained by varying the common signal at fixed rare signal.

\paragraph{Noise, budget, and finite prevalence.}

The comparison need not be monotone in noise: with positive common signal, U
wins at sufficiently small and sufficiently large noise, while SD can win
between them. For fixed $r,\nu>0$ and finite $s\geq0$, the limiting risks
also satisfy $L_S(\tau)/L_U(\tau)\to4/5$ as $\tau\to\infty$. Thus SD
eventually wins as the rescaled budget grows, although convergence can be slow.
This is an ordered limit, taking $q\to0$ first; it makes no claim about a joint
limit with $M_qb\to\infty$. The expansions and examples are in
\cref{sec:oracle-noise-budget}.

Exact-moment searches at six prevalences recover the predicted signs away from
limiting ties at the smallest prevalence. At larger prevalence, a high-noise
example instead favors SD on the finite grids where the limit favors U.
\Cref{sec:oracle-finite} reports this disagreement and the search refinements.
These numerical comparisons assess the approximation and do not certify
continuous finite-prevalence optima.
